\documentclass[
  spanish,
  aspectratio=169,
  xcolor={dvipsnames,table}
]{beamer}

\geometry{paperwidth=19.2cm,paperheight=10.8cm}
\usepackage[utf8]{inputenc}
\usepackage[T1]{fontenc}
\usepackage[spanish,es-tabla]{babel}
\usepackage[scaled=.96]{helvet}
\renewcommand{\familydefault}{\sfdefault}
\usepackage{booktabs}
\usepackage{graphicx}
\usepackage{ragged2e}
\usepackage{tikz}
\usepackage{hyperref}
\usetikzlibrary{calc,positioning,fit,arrows.meta,shadows.blur}

\newcommand{\studentname}{Martin Jonathan de la Cruz}
\newcommand{\studentshort}{M. J. de la Cruz}
\newcommand{\studentid}{ES2611202040}
\newcommand{\universityname}{Universidad Abierta y a Distancia de México}
\newcommand{\facultyname}{Licenciatura en Derecho}
\newcommand{\coursename}{Garantías Constitucionales}
\newcommand{\coursecode}{LDE-S2B1}
\newcommand{\activitytitle}{Mapa conceptual}
\newcommand{\activitysubtitle}{Marco teórico de los derechos humanos y sus principios}
\newcommand{\documentsubject}{Actividad 3 - Garantías Constitucionales}
\newcommand{\activityweek}{Semanas 2 y 3}
\newcommand{\teachingfigure}{Aarón López Pérez}
\newcommand{\deliverydate}{\today}
\newcommand{\locationname}{Roma Norte, Ciudad de México}
\newcommand{\departmentlogo}{img/departamentos/UnADM.pdf}

\definecolor{unadmGreenDark}{HTML}{174A3A}
\definecolor{unadmGreen}{HTML}{5F8F3A}
\definecolor{unadmGold}{HTML}{B88A2A}
\definecolor{unadmPaper}{HTML}{F6F7F2}
\definecolor{unadmInk}{HTML}{1F2A24}
\definecolor{mcLeaf}{HTML}{EAF1E2}

\mode<presentation>{
  \usetheme{default}
  \usefonttheme{professionalfonts}
  \setbeamertemplate{navigation symbols}{}
  \setbeamertemplate{blocks}[rounded][shadow=false]
}
\setbeamersize{text margin left=0.60cm,text margin right=0.60cm}
\setbeamercolor{background canvas}{bg=white}
\setbeamercolor{normal text}{fg=unadmInk,bg=white}
\setbeamercolor{structure}{fg=unadmGreenDark}
\setbeamercolor{frametitle}{fg=white,bg=unadmGreenDark}
\setbeamercolor{block title}{fg=white,bg=unadmGreen}
\setbeamercolor{block body}{fg=unadmInk,bg=unadmPaper}
\setbeamerfont{title}{size=\LARGE,series=\bfseries}
\setbeamerfont{frametitle}{size=\large,series=\bfseries}
\setbeamertemplate{itemize item}{\textcolor{unadmGreen}{\large$\blacktriangleright$}}

\setbeamertemplate{frametitle}{
  \nointerlineskip
  \begin{beamercolorbox}[wd=\textwidth,ht=1.02cm,dp=0.22cm,leftskip=0.35cm,rightskip=0.25cm]{frametitle}
    \insertframetitle\hfill\includegraphics[height=0.60cm]{\departmentlogo}
  \end{beamercolorbox}
  {\color{unadmGold}\rule{\textwidth}{1.3pt}}
}

\setbeamertemplate{footline}{
  \leavevmode
  \hbox{
    \begin{beamercolorbox}[wd=.34\paperwidth,ht=2.7ex,dp=1ex,leftskip=1em]{author in head/foot}\color{white}\studentshort\end{beamercolorbox}
    \begin{beamercolorbox}[wd=.40\paperwidth,ht=2.7ex,dp=1ex,center]{title in head/foot}\color{white}\coursename\end{beamercolorbox}
    \begin{beamercolorbox}[wd=.26\paperwidth,ht=2.7ex,dp=1ex,rightskip=1em plus 1fil]{date in head/foot}\color{white}\coursecode\hfill\insertframenumber/\inserttotalframenumber\end{beamercolorbox}
  }
}
\setbeamercolor{author in head/foot}{bg=unadmGreenDark}
\setbeamercolor{title in head/foot}{bg=unadmGreen}
\setbeamercolor{date in head/foot}{bg=unadmGreenDark}

\title[\coursecode]{\activitytitle}
\subtitle{\activitysubtitle}
\author[\studentshort]{\studentname\\\footnotesize Matrícula: \studentid}
\institute[UnADM]{\universityname\\\facultyname}
\date[\activityweek]{\locationname\\\activityweek\\\deliverydate}

\tikzset{
  pcroot/.style={rectangle, rounded corners=3pt, fill=unadmGreenDark, text=white,
    font=\bfseries\scriptsize, align=center, inner sep=4pt, text width=2.9cm},
  pcbranch/.style={rectangle, rounded corners=2pt, fill=unadmGreen, text=white,
    font=\bfseries\tiny, align=center, inner sep=3pt, text width=2.2cm},
  pclink/.style={-{Latex[length=1.4mm]}, draw=unadmGreenDark!70, line width=0.6pt}
}

\begin{document}

\begin{frame}[plain]
  \begin{tikzpicture}[remember picture,overlay]
    \fill[unadmGreenDark] (current page.south west) rectangle ([xshift=0.58\paperwidth]current page.north west);
    \fill[unadmGreen] ([xshift=0.58\paperwidth]current page.south west) rectangle (current page.north east);
    \node[anchor=center,opacity=0.10] at (current page.center) {\includegraphics[width=0.72\paperwidth]{\departmentlogo}};
    \node[anchor=north west] at ([xshift=0.58cm,yshift=-0.46cm]current page.north west) {\includegraphics[width=2.55cm]{\departmentlogo}};
  \end{tikzpicture}
  \vspace*{1.62cm}
  {\color{white}\fontsize{24}{29}\selectfont\bfseries \activitytitle\par}
  \vspace{0.28cm}
  {\color{white!86}\large \activitysubtitle}\\[0.16cm]
  {\color{white!82}\normalsize \documentsubject}\\[0.35cm]
  {\color{unadmGold}\rule{0.58\linewidth}{1.3pt}}\\[0.35cm]
  {\color{white}\studentname}\\
  {\color{white!82}\footnotesize \facultyname}
\end{frame}

\begin{frame}{Propósito de la actividad}
  \begin{block}{Mapa conceptual del Tema 1}
    Ubicar los antecedentes históricos y reconocer el concepto de derechos humanos y los principios dogmáticos que los rigen.
  \end{block}
  \begin{itemize}
    \item Actividad sumativa del bloque (semanas 2 y 3).
    \item Síntesis jerárquica del marco teórico de los derechos humanos.
    \item Base conceptual para el estudio de los medios de control constitucional.
  \end{itemize}
\end{frame}

\begin{frame}{Estructura del mapa conceptual}
  \begin{center}
  \begin{tikzpicture}[node distance=0.35cm and 0.35cm]
    \node[pcroot] (r) {Derechos humanos y sus principios};
    \node[pcbranch, below left=1.1cm and 2.7cm of r] (h) {1.1 Antecedentes históricos};
    \node[pcbranch, below left=1.1cm and 0.5cm of r] (g) {1.2 Generaciones};
    \node[pcbranch, below=1.2cm of r] (c) {1.3 Concepto};
    \node[pcbranch, below right=1.1cm and 0.5cm of r] (p) {1.4 Principios (art. 1.º)};
    \node[pcbranch, below right=1.1cm and 2.7cm of r] (ga) {1.5 Garantía};
    \draw[pclink] (r) -- (h);
    \draw[pclink] (r) -- (g);
    \draw[pclink] (r) -- (c);
    \draw[pclink] (r) -- (p);
    \draw[pclink] (r) -- (ga);
  \end{tikzpicture}
  \end{center}
  {\footnotesize Cinco ejes jerarquizados desde un concepto raíz común.}
\end{frame}

\begin{frame}{Antecedentes y generaciones}
  \begin{columns}[T]
    \begin{column}{0.48\linewidth}
      \begin{block}{Antecedentes históricos}
        Proceso de luchas y documentos que reconocen la dignidad y limitan el poder.
      \end{block}
    \end{column}
    \begin{column}{0.48\linewidth}
      \begin{block}{Generaciones}
        1.ª civiles y políticos; 2.ª económicos, sociales y culturales; 3.ª colectivos.
      \end{block}
    \end{column}
  \end{columns}
  \vspace{0.2cm}
  Las generaciones ordenan su desarrollo, sin jerarquía entre ellas.
\end{frame}

\begin{frame}{Concepto y principios}
  \begin{block}{Derecho humano}
    Prerrogativa ligada a la dignidad; obliga al Estado a respetar, proteger, promover y garantizar (art. 1.º constitucional).
  \end{block}
  \begin{itemize}
    \item \textbf{Universalidad} e \textbf{interdependencia}.
    \item \textbf{Indivisibilidad}.
    \item \textbf{Progresividad} y \textbf{no regresión}.
  \end{itemize}
\end{frame}

\begin{frame}{Concepto de garantía}
  \begin{columns}[T]
    \begin{column}{0.48\linewidth}
      \begin{block}{Derecho}
        El bien o facultad protegida.
      \end{block}
    \end{column}
    \begin{column}{0.48\linewidth}
      \begin{block}{Garantía}
        Mecanismo jurídico que lo hace exigible.
      \end{block}
    \end{column}
  \end{columns}
  \vspace{0.25cm}
  Sin garantías eficaces, los derechos quedan como declaraciones formales.
\end{frame}

\begin{frame}{Cierre}
  \begin{block}{Aprendizaje principal}
    El mapa articula historia, generaciones, concepto, principios y garantías como una secuencia lógica y jerarquizada.
  \end{block}
  \begin{itemize}
    \item El eje articulador es la distinción entre derecho y garantía.
    \item Prepara el estudio de los medios de control constitucional.
    \item Fuentes: Álvarez Icaza, Rabinovich Berkman y Marcano Salazar.
  \end{itemize}
\end{frame}

\end{document}
